\section{Der Jahresabschluss 2025}
\label{sec:abschluss}

\subsection{\"Uberblick}

\ueberblicktabelle

Umsatz und EBITDA vor Restrukturierung erreichten \Mio{\Umsatz}{EUR} und
\Mio{\EbitdaVr}{EUR}. Beide Gr\"o{\ss}en wuchsen um etwa denselben absoluten Betrag
(Tabelle~\ref{tab:ueberblick}); das Wachstum kam 2025 nicht aus der Menge, sondern aus
der Marge.

\subsection{Mehrjahresreihe}

\reihentabelle

\subsection{Auftragseingang als vorlaufende Gr\"o{\ss}e}

Der Auftragseingang \"ubersteigt den Umsatz in jedem Jahr der Reihe. 2025 betrug er
\Mio{\Auftragseingang}{EUR}\Q{GeaGB}{2}, der Auftragsbestand zum Jahresende
\Mio{\Auftragsbestand}{EUR}\Q{GeaGB}{44}; gemessen am Umsatz des Jahres entspricht der
Bestand \Auftragsreichweite{} Monaten Arbeit.

Im ersten Halbjahr 2026 stieg der Auftragseingang auf \Mio{\HjAuftragseingang}{EUR} nach
\Mio{\HjAuftragseingangVorjahr}{EUR} im Vorjahreszeitraum, ein Zuwachs von
\HjAuftragWachstum; der Auftragsbestand erreichte \Mio{\HjAuftragsbestand}{EUR} und liegt
damit \"uber dem Stand zum Jahresende\Q{GeaHJ}{2}. Das EBITDA vor Restrukturierung des
Halbjahres betrug \Mio{\HjEbitdaVr}{EUR} nach \Mio{\HjEbitdaVrVorjahr}{EUR}, die Marge
\HjMarge{} gegen\"uber \Pz{\EbitdaMarge} im Gesamtjahr 2025.

\bk{Der Auftragseingang ist bei einem Anlagenbauer die einzige Gr\"o{\ss}e, die dem Umsatz
zeitlich vorausgeht, und er ist im Halbjahr st\"arker gestiegen als in jedem der vier
Jahre davor. Zusammen mit einem Bestand, der \"uber dem Jahresendstand liegt, ist der
Umsatz der kommenden vier bis f\"unf Quartale weitgehend vorgezeichnet. Das ist der
belastbarste Grund, der f\"ur GEA spricht, und Abschnitt~\ref{sec:kurs} pr\"uft, ob der
Kurs ihm folgt.}

\subsection{Bilanz und Finanzlage}

Die Eigenkapitalquote betrug zum Jahresende \EkQuote, das Eigenkapital
\Mio{\Eigenkapital}{EUR}\Q{GeaGB}{312}. Statt Nettoverschuldung weist GEA eine
Nettoliquidit\"at von \Mio{\Nettoliquiditaet}{EUR} aus\Q{GeaGB}{2}. Der ROCE erreichte
\Pz{\Roce}.

\bk{Ein Unternehmen ohne Nettoschulden hat keine Finanzierungsfrage, sondern eine
Verwendungsfrage. Das ist f\"ur die Flow-Analyse in Abschnitt~\ref{sec:flow} entscheidend:
Es gibt keinen Zins, der einen Teil des Zuflusses vorab bindet, und keine F\"alligkeit, die
eine Investition erzwingt. Der ROCE von \Pz{\Roce} ist die Kehrseite derselben
Tatsache: Das eingesetzte Kapital ist klein, weil der Anlagenbau beim Kunden montiert und
nicht auf eigenen Anlagen produziert.}

\subsection{Ein Widerspruch im Bericht}

Der Prognose-Ist-Vergleich auf Seite~40 nennt als Ergebnis 2025 eine EBITDA-Marge vor
Restrukturierung von \Pz{\MargeVergleich}\Q{GeaGB}{40}. Die Kennzahlenseite nennt
\Pz{\EbitdaMarge}\Q{GeaGB}{2}. Aus den ebenfalls gedruckten Werten
\Mio{\EbitdaVr}{EUR} und \Mio{\Umsatz}{EUR} ergibt sich \MargeBerechnet. Beide Zahlen
stehen im selben Bericht; die Rechnung st\"utzt die Kennzahlenseite.

\bk{Die \Pz{\MargeVergleich} sind zugleich die Untergrenze der zuletzt angepassten
Prognose (Tabelle~\ref{tab:prognose}). Der wahrscheinlichste Hergang ist, dass in die
Ergebnisspalte die Prognose statt des Ergebnisses geraten ist. Der Bericht rechnet
durchg\"angig mit \Pz{\EbitdaMarge}, weil diese Zahl aus zwei anderen gedruckten Zahlen
folgt. Festzuhalten bleibt, dass eine der beiden Angaben im gepr\"uften Gesch\"aftsbericht
falsch ist.}

\subsection{Was mit dem Geld geschah}

Der freie Zufluss betrug 2025 \Mio{\FcfVoll}{EUR}. Davon flossen
\Mio{\DividendeSumme}{EUR} als Dividende und \Mio{\RueckkaufSumme}{EUR} in
Aktienr\"uckk\"aufe ab, zusammen \AusschuettungAnteil{} des Zuflusses.

\dividendentabelle

\bk{Die zur\"uckgerechnete Dividende je Aktie best\"atigt die auf Seite~16 beschriebene
Politik aus den Zahlungsstr\"omen heraus: bis 2024 Schritte von rund f\"unf Cent, danach
gr\"o{\ss}ere. Die Reihe zeigt zugleich, was die R\"uckk\"aufe bewirken: Die Zahlung stieg
2025 st\"arker als die Dividende je Aktie, weil weniger Aktien mehr je St\"uck erhalten.}
